% SPDX-License-Identifier: Apache-2.0
% covers: BeatMaster
\datasheet{BeatMaster}{An AXI4 host that performs one beat at a time as a burst of one, and answers each, repeating the last answer to an ask that repeats its tag.}

\dsfacts{\code{lib/parts/src/remote/stream.rs}}{\code{txhdl\_parts::remote::stream::BeatMaster}}%
{\code{//docs:examples}}

\subsection*{Function}

\code{BeatMaster} takes a \code{BeatAsk} from \code{StreamLink}, a
program's read or write of one 16-byte beat, performs it on an AXI4
port of 128-bit data, 64-bit addresses and 4-bit identifiers as a burst
of one beat, and sends the \code{BeatAnswer} back. It is how FireSim's
host driver reaches its CPU-managed streams over Ethernet (issue 1984).

A stream read takes data out of a queue and a stream write puts it in,
so neither may be performed twice. As with \code{LiteMaster}, an ask
that repeats the tag of the last transfer is answered from what was
kept of it, the beat included, and not performed again.

A program starts with a tag it cannot know the board has not kept from
an earlier run, and no stream transfer is harmless to repeat or to skip.
So its first ask is a no-op: \code{BeatMaster} answers it without
touching the port and keeps its tag as the last one, and every tag the
program uses after it differs from the one kept.

\subsection*{Parameters}

None. The port is AXI4 of 64-bit addresses, 128-bit data, a 16-bit
strobe and 4-bit identifiers.

\dstables{BeatMaster}

\subsection*{Behaviour}

\begin{itemize}
\item An ask is taken when nothing is on the port and no answer is
waiting to go.
\item An ask whose tag is the last transfer's is answered again from
\code{last\_tag}, \code{last\_data} and \code{last\_err}, and not
performed. A no-op is never taken as a repeat, so it cannot bring back
a kept error at the start of a run.
\item A no-op is answered without the port: it sets the kept tag, a
zero beat and no error.
\item A write sends one address phase, identifier 0, length 0, size 16
bytes, incrementing, and one beat marked last; a read sends the address
phase on \code{ar}.
\item The response is taken only once the request has gone whole. A
response other than \code{Okay} becomes an answer with its error bit
set; a write answers a zero beat.
\end{itemize}

\subsection*{Verification}

\begin{itemize}
\item The tests in \code{lib/parts/src/remote/stream.rs}, in
\code{//lib/parts:parts\_test}:
\code{a\_write\_reaches\_the\_port\_whole}, one beat of 16 bytes;
\code{reads\_take\_the\_queue\_in\_order};
\code{a\_repeated\_read\_is\_answered\_again\_and\_not\_performed\_again}
and \code{a\_repeated\_write\_is\_not\_performed\_again};
\code{a\_noop\_is\_answered\_and\_never\_reaches\_the\_port}; and\newline
\code{a\_noop\_on\_a\_failed\_transfers\_tag\_does\_not\_bring\_back\_the\_error}.
\item \code{//lib/examples:ex\_stream}, and the build simulates its
netlist against that run under nvc and Verilator:\newline
\code{//docs:streamlink\_sim\_beat\_master\_beat\_master\_tb\_test}\newline
\code{//docs:streamlink\_vsim\_beat\_master\_test}.
\end{itemize}

\subsection*{Used by}

\code{StreamBridge}, which joins it to \code{StreamLink}; \code{ex\_stream};
and \code{//firesim/ax7a200b} through \code{StreamBridge<5>}.

\subsection*{Limits}

One beat at a time, a burst of one. A repeat is known by its tag alone,
so each new ask must change its tag. Only the last answer is kept.
